%% ma: L12-B06-0023
%% lop: 12
%% chuong: 2
%% bai: 6
%% dang: 12.6.15
%% loai: TN4
%% muc: TH
%% dap_an: "B"
%% nguon: "LaTeX giáo viên gửi 21/10/2026 (M05, Vectơ và hệ trục tọa độ trong không gian, 107 câu), Phần I Câu 32"
%% kien_thuc: "Bài 6 (tích vô hướng, ứng dụng tính công)"
%% trang_thai: chinh-thuc
%% ghi_chu: "Giáo viên đã duyệt 21/10/2026."
\begin{ex}
Người ta dùng một lực kéo $\overrightarrow{F}$ có phương hợp với mặt phẳng ngang một góc $30^\circ$ và có độ lớn $200$ N để kéo một kiện hàng trượt thẳng trên mặt sàn nằm ngang một quãng đường dài $5$ m. Công sinh ra bởi lực $\overrightarrow{F}$ là
\choice{$500$ J}{\True $500\sqrt3$ J}{$1000$ J}{$1000\sqrt3$ J}
\loigiai{Gọi $\vec{d}$ là vectơ độ dời, $|\vec{d}|=5$ và $(\overrightarrow{F},\vec{d})=30^\circ$. Công $A=\overrightarrow{F}\cdot\vec{d}=200\cdot5\cdot\cos30^\circ=500\sqrt3$ (J).}
\end{ex}
